\section{Research Gap and Positioning of Motivo}\label{sec:motivo-position}

The systems discussed in this chapter illustrate a wide spectrum of approaches to music programming and representation,
each addressing different aspects of musical creation.
High-level composition systems such as \textbf{Sonic Pi}, \textbf{Tidal}, and \textbf{HMusic} prioritize musical
abstraction, expressiveness, and interactive composition workflows.
However, they often rely on sequential or pattern-based models that can limit explicit control over complex temporal
relationships.

Systems for real-time audio control and synthesis, including \textbf{ChucK} and \textbf{SuperCollider}, provide precise
control over timing and sound synthesis.
While these systems offer high flexibility, they operate at a relatively low level of abstraction, requiring users to
manage details related to signal processing, concurrency, and timing explicitly.

Formats for musical representation and exchange such as \textbf{MIDI} and \textbf{ABC notation} focus on encoding
musical information for storage, transmission, and playback\@.
While widely adopted, these formats serve different goals from a composition-oriented domain-specific
language\@: \textbf{MIDI} is effective
as an interchange and playback format, and \textbf{ABC notation} is effective as a human-readable notation format.
Neither is intended to serve as a high-level, executable model for structured composition.

Table~\ref{tab:related-work-comparison} summarizes the main systems and languages discussed in this chapter and
highlights the design trade-offs that motivate Motivo\@.

\begin{table}[htbp]
  \centering
  \footnotesize
  \renewcommand{\arraystretch}{1.5}
  \begin{tabular}{p{2.2cm}p{2.5cm}p{4cm}p{5cm}}
    \toprule
    \textbf{System / Language} & \textbf{Class} & \textbf{Temporal / structural model} & \textbf{Main gap
    relative to this thesis} \\
    \midrule
    Sonic Pi & High-level composition & Sequential execution with notes, sleeps, and simple patterns &
    Limited structural abstraction and weak support for complex temporal relationships \\
    Tidal & High-level composition & Time-varying patterns with functional composition over discrete or
    continuous time & Haskell-centric syntax and focus on event streaming rather than direct compilation or
    explicit timeline structure \\
    HMusic & High-level composition & Grid-like patterns and parallel tracks for rhythmic structure & Relies
    on Sonic Pi, remains pattern-centric, and does not model higher-level timelines explicitly \\
    ChucK & Real-time audio control and synthesis & Explicit time advancement and synchronized concurrent
    shreds & Low-level time management makes high-level compositional structure harder to express \\
    Super\-Collider & Real-time audio control and synthesis & Object-oriented and functional model with
    UGens, patterns, and streams & Powerful but synthesis-oriented, with substantial complexity and weaker
    support for structured high-level composition \\
    MIDI & Musical representation & Discrete message stream of note and control events & Strong backend and
    interchange format, but too low-level to serve as the primary model for structured composition \\
    ABC notation & Musical representation & Declarative textual notation with implicit ordering, bars, and
    voices & Clear notation format, but not an executable high-level composition language with explicit
    temporal abstraction \\
    \bottomrule
  \end{tabular}
  \caption{Comparison of the main systems and languages discussed in this
  chapter.}\label{tab:related-work-comparison}
\end{table}

From this analysis, several gaps relative to the goals of this thesis emerge.
First, many systems lack an explicit and flexible model for representing time, often relying on implicit sequencing or
pattern repetition.
Second, there is a gap between low-level audio programming languages and high-level compositional tools, with few
systems providing both expressive power and structured abstraction.
Third, existing approaches often conflate composition with execution, making it difficult to separate musical intent
from performance details.

These gaps motivate the development of Motivo, which aims to
provide a timeline-based model for music composition, support deterministic and structured representations of musical
ideas, and enable modular and extensible compilation to multiple output formats.

The next chapter moves from the surrounding ecosystem to the design of Motivo itself. It explains the
language model, the compiler architecture, and the lowering process that turns recursive musical structure
into MIDI events.
